\documentclass[11pt,a4paper]{exam}
\noprintanswers
%\printanswers

\usepackage[T1]{fontenc}
\usepackage{lmodern}
\usepackage{amsmath,amssymb}
\usepackage[margin=2cm,includeheadfoot,headheight=15pt,headsep=18pt]{geometry}
\usepackage{needspace}
\usepackage[hidelinks]{hyperref}

\newcommand{\CourseTitle}{Topics in Math. Sci. VII 2026}
\newcommand{\AssignmentNumber}{1}
\newcommand{\Deadline}{15th November 2026}

\pagestyle{headandfoot}
\header{}{\small\CourseTitle}
       {\small Assignment \AssignmentNumber}
\headrule
\footer{}{\small Submit via email.  Deadline: \Deadline}{\small\thepage}

\newcommand{\kk}{\Bbbk}
\newcommand{\bbD}{\mathbb D}
\newcommand{\grstar}{(*)}
\DeclareMathOperator{\Hom}{Hom}
\DeclareMathOperator{\End}{End}
\DeclareMathOperator{\rad}{rad}
\newcommand{\SymAlg}{\mathsf{S}}
\newcommand{\op}{\mathrm{op}}
\newcommand{\id}{\mathrm{id}}

\qformat{\textbf{Exercise \thequestion.\ \thequestiontitle}\hfill}
\renewcommand{\questionshook}{%
  \setlength{\leftmargin}{0pt}%
  \setlength{\labelwidth}{-\labelsep}%
  \setlength{\itemsep}{13pt}%
}
\renewcommand{\thepartno}{\arabic{partno}}
\renewcommand{\partlabel}{(\thepartno)}
\renewcommand{\partshook}{\setlength{\itemsep}{4pt}}
\setlength{\parindent}{0pt}
\setlength{\parskip}{4pt}
\setlength{\emergencystretch}{1em}
\raggedbottom

\newcommand{\comm}[1]{}


\begin{document}

Complete at least 2 exercises.

Throughout, $\kk$ is a field, and all modules are right modules unless a
left action is specified.


\begin{questions}

\titledquestion{}

\begin{parts}
\part Let $A=\bigoplus_{d\geq0}A_d$ be a graded algebra with $S=A_0$ finite-dimensional semisimple.
\begin{subparts}
\subpart Prove that
$A/A_+\cong S$ and $\rad A\subset A_+$.
\begin{solution}
The projection $A\to A_0$ onto degree $0$ is an algebra homomorphism with kernel $A_+$, so $A/A_+\cong S$.
For a maximal right ideal $\mathfrak m\subset S$, its preimage $\mathfrak m+A_+$ is a right ideal with $A/(\mathfrak m+A_+)\cong S/\mathfrak m$ simple, hence maximal. Therefore
\[
  \rad A\ \subset\bigcap_{\mathfrak m}(\mathfrak m+A_+)=\Bigl(\bigcap_{\mathfrak m}\mathfrak m\Bigr)+A_+=\rad S+A_+=A_+,
\]
since $\rad S=0$ for semisimple $S$.
\end{solution}
\subpart Prove that if $A_+^N=0$ for some $N\geq1$ (for example if $\dim_\kk A<\infty$), then $\rad A=A_+$.

\begin{solution}
Let $L$ be a simple right $A$-module. Then $LA_+$ is a submodule, so $LA_+\in\{0,L\}$, and $LA_+=L$ would give $L=LA_+^N=0$. Hence $LA_+=0$.
Applying this to $L=A/\mathfrak m$ for a maximal right ideal $\mathfrak m$ gives $A_+=1\cdot A_+\subset\mathfrak m$. So $A_+\subset\rad A$, and equality follows from (i).
\end{solution}
\subpart Show that for $A=\kk[x]$, $\deg x=1$, one has $\rad A=0\neq A_+$.
\begin{solution}
The maximal ideals of $\kk[x]$ are $(p)$ with $p$ irreducible, and there are infinitely many irreducible polynomials (if $p_1,\dots,p_m$ were all of them, any irreducible factor of $p_1\cdots p_m+1$ would be a new one).
A non-zero polynomial has finitely many irreducible factors, so $\bigcap_p(p)=0$, while $A_+=(x)\neq0$.
\end{solution}
\subpart Suppose that $\dim_\kk A<\infty$ and that $A$ is generated in degree $1$ over $S$. Prove that
\[
  \rad^kA=\bigoplus_{d\geq k}A_d,\qquad\rad^kA/\rad^{k+1}A\cong A_k\qquad(k\geq0).
\]
\begin{solution}
By (ii), $\rad A=A_+$, so $\rad^kA=A_+^k$. Clearly $A_+^k\subset\bigoplus_{d\geq k}A_d$. Conversely, generation in degree $1$ gives $A_d=A_1^d\subset A_+^d\subset A_+^k$ for $d\geq k$.
Hence $\rad^kA/\rad^{k+1}A=\bigoplus_{d\geq k}A_d\big/\bigoplus_{d\geq k+1}A_d\cong A_k$.
\end{solution}
\end{subparts}

\part Let $A=\kk A_n$, $P_i:=e_{ii}A$ with the inherited grading, and
  $L_j:=e_{jj}A/e_{jj}A_+$ as graded module concentrated in degree 0.
\begin{subparts}
\subpart Prove that for $k\geq0$, $\rad^kP_i=\operatorname{span}_\kk\{e_{im}\mid m\geq i+k\}$.
\begin{solution}
By $\rad M=M\rad A$ and (1)(iv), $\rad^kP_i=P_i(\rad A)^k=P_iA_{\geq k}$. Now $e_{ip}e_{pq}=e_{iq}$ with $p\geq i$ and $q-p\geq k$ ranges over all $e_{iq}$ with $q\geq i+k$ (take $p=i$ for $q\geq i+k$), and every product $e_{ip}A_{\geq k}$ lies in this span.
\end{solution}
\subpart Prove that for $0\leq k\leq n-i$,
$\rad^kP_i/\rad^{k+1}P_i\ \cong\ L_{i+k}(-k)$ as graded right $A$-modules, and deduce that
\[
  P_i\cong P_j(d)\ \text{ as graded modules}\quad\Longleftrightarrow\quad(i,d)=(j,0).
\]
\begin{solution}
By (i) the quotient has basis the class $\bar e$ of $e_{i,i+k}$, of degree $k$. For a matrix unit $e_{pq}$ we have $\bar e\,e_{pq}=\delta_{i+k,p}\,\overline{e_{iq}}$, which vanishes in the quotient unless $q=i+k$. So $e_{i+k,i+k}$ acts as the identity and all other matrix units act by $0$:
the module is $L_{i+k}$ with its generator placed in degree $k$, i.e.\ $L_{i+k}(-k)$.
A graded isomorphism $P_i\cong P_j(d)$ induces an isomorphism of tops $L_i=P_i/\rad P_i\cong(P_j/\rad P_j)(d)=L_j(d)$. The left side is concentrated in degree $0$ and the right side in degree $-d$, so $d=0$; and $e_{ii}$ acts non-trivially on the left, so $i=j$. The converse is clear.
\end{solution}
\end{subparts}
\end{parts}

\titledquestion{}
\begin{parts}
\part Let $A=\Bbbk A_3$, $M=e_{22}A$, and
consider $\bbD(M)_d=\Hom_\kk(M_{-d},\kk)$ with $(af)(m)=f(ma)$.
\begin{subparts}
\subpart Find a homogeneous basis of $\bbD(M)$, its degrees, and all non-zero
actions of matrix units on it.
\begin{solution}
$M=\kk e_{22}\oplus\kk e_{23}$ with degrees $0,1$. Let $\varepsilon_{22},\varepsilon_{23}$ be the dual basis;
$\deg\varepsilon_{22}=0$, $\deg\varepsilon_{23}=-1$. Since $(af)(m)=f(ma)$ and
$e_{ij}e_{pq}=\delta_{jp}e_{iq}$, the only non-zero actions are
$e_{22}\varepsilon_{22}=\varepsilon_{22}$, $e_{33}\varepsilon_{23}=\varepsilon_{23}$,
$e_{23}\varepsilon_{23}=\varepsilon_{22}$.
(E.g.\ $(e_{23}\varepsilon_{23})(e_{22})=\varepsilon_{23}(e_{22}e_{23})=1$, while $e_{11},e_{12},e_{13}$
act by $0$ because $e_{22}e_{1q}=e_{23}e_{1q}=0$.)
\end{solution}
\subpart Show that $\Bbbk e_{13}\subset Ae_{33}$ is a graded left submodule and
that $\bbD(M)\cong(Ae_{33}/\Bbbk e_{13})(1)$ as graded left $A$-modules.
\begin{solution}
$Ae_{33}=\operatorname{span}\{e_{13},e_{23},e_{33}\}$ and $e_{ij}e_{13}=\delta_{j1}e_{i3}$, so
$\Bbbk e_{13}$ is a (homogeneous) submodule. In $N=Ae_{33}/\Bbbk e_{13}$ the classes $\bar e_{33},\bar e_{23}$
have degrees $0,1$ and the non-zero actions are $e_{33}\bar e_{33}=\bar e_{33}$,
$e_{23}\bar e_{33}=\bar e_{23}$, $e_{22}\bar e_{23}=\bar e_{23}$. In $N(1)$ the degrees are $-1,0$.
The linear map $\varepsilon_{23}\mapsto\bar e_{33}$, $\varepsilon_{22}\mapsto\bar e_{23}$ is bijective, preserves
degrees, and matches the non-zero actions of part (i); as matrix units span $A$, it is $A$-linear.
\end{solution}
\subpart Deduce that $\bbD(M)$ is not isomorphic to any
$(Ae_{jj})(r)$, $j\in\{1,2,3\}$, $r\in\mathbb Z$.
\begin{solution}
$\dim\bbD(M)=2$ and $\dim Ae_{jj}=j$, so only $j=2$ is possible. But $e_{33}$ acts by $0$ on $Ae_{22}=\operatorname{span}\{e_{12},e_{22}\}$
while $e_{33}\varepsilon_{23}=\varepsilon_{23}\neq0$; an isomorphism commutes with $e_{33}$.
\end{solution}
\end{subparts}

\part Let $A=\kk[x,y]/(x^2,y^3)$ with $\deg x=\deg y=1$, and let $V=A_1$.
\begin{subparts}
\subpart Compute the Hilbert series $h_A(t)$.
\begin{solution}
$A\cong\kk[x]/(x^2)\otimes_\kk\kk[y]/(y^3)$ has basis $x^ay^b$ with $a<2$, $b<3$, so
\[
  h_A(t)=(1+t)(1+t+t^2)=1+2t+2t^2+t^3.
\]
\end{solution}
\subpart Let $R=\ker\bigl(V\otimes_\kk V\to A_2\bigr)$ and $Q=T_\kk(V)/(R)$. Show that
\[
  Q\cong\kk[x,y]/(x^2),\qquad h_Q(t)=\frac{1+t}{1-t}.
\]
\begin{solution}
$A_2$ has basis $xy,y^2$, so $\dim R=4-2=2$ and $R=\operatorname{span}\{x\otimes x,\ x\otimes y-y\otimes x\}$.
Thus $Q=\kk\langle x,y\rangle/(x^2,xy-yx)=\kk[x,y]/(x^2)$, with basis $x^ay^b$, $a\in\{0,1\}$, $b\geq0$; this has $\dim Q_d=2$ for $d\geq1$, so $h_Q=1+\sum_{d\geq1}2t^d=\frac{1+t}{1-t}$.
\end{solution}
\subpart Prove that $A$ is not quadratic, and find the smallest degree in which $A$ and $Q$ differ.
\begin{solution}
$A$ is generated in degree $1$, and any quadratic presentation $A\cong T_\kk(V)/(R')$ has $V=A_1$ and $R'=R$ (Lecture~2), hence would give $A\cong Q$.
The natural surjection $Q\to A$ is an isomorphism in degrees $\leq2$, but $\dim Q_3=2>1=\dim A_3$, with $y^3\neq0$ in $Q$. So $A\not\cong Q$ and $A$ is not quadratic; the smallest degree of difference is $3$.
\end{solution}
\end{subparts}


\end{parts}

\titledquestion{}
\begin{parts}
\part For $n\geq3$ and $Z\subset\{1,\dots,n-2\}$ let
\[
  A_Z=\kk A_n/(e_{k,k+2}\mid k\in Z)
\]
with the quotient grading.
\begin{subparts}
\subpart Show that
\[
  \{e_{ii}\}_{i}\ \cup\ \{e_{ij}\mid i<j,\ \nexists\,k\in Z:\ i\leq k\leq j-2\}
\]
is a basis of $A_Z$.
\begin{solution}
Let $I=\operatorname{span}\{e_{ij}\mid\exists k\in Z:\ i\leq k,\ k+2\leq j\}$. Since $e_{ij}=e_{ik}e_{k,k+2}e_{k+2,j}$, $I$ lies in $(e_{k,k+2}:k\in Z)$;
conversely $I$ is a two-sided ideal containing the $e_{k,k+2}$ (a product $e_{pq}e_{ij}e_{rs}\neq0$ has $p\leq i\leq k$ and $k+2\leq j\leq s$).
So $I=(e_{k,k+2}:k\in Z)$ and the remaining matrix units form a basis of the quotient.
\end{solution}
\subpart Write $A_Z=T_S(V)/(R)$ with explicit $S$, $V$, and $R\subset V\otimes_SV$.
\begin{solution}
\[
  S=\bigoplus_{i=1}^n\kk e_{ii},\qquad V=\bigoplus_{i<n}\kk e_{i,i+1},\qquad R=\operatorname{span}\{e_{k,k+1}\otimes e_{k+1,k+2}\mid k\in Z\}.
\]
Indeed $T_S(V)\cong\kk A_n$ (Lecture 2), and under this
isomorphism $(R)$ goes to the ideal generated by the products $e_{k,k+1}e_{k+1,k+2}=e_{k,k+2}$.
\end{solution}
\end{subparts}

\part Let $S=\kk e_1\oplus\kk e_2$ and $V=\kk a\oplus\kk b$ with
\[
  a\in e_1Ve_2,\qquad b\in e_2Ve_1,
\]
and $e_iVe_j=0$ otherwise.
\begin{subparts}
\subpart Find a basis of $T_S(V)_d$ for all $d\geq0$, and compute $h_{T_S(V)}(t)$.
\begin{solution}
Positive-degree basis words of $T_S(V)$ are the composable (alternating) words. For each $d\geq1$ there are exactly two:
$aba\cdots$ ($d$ letters, from $e_1$) and $bab\cdots$ (from $e_2$); for $d=0$ the basis is $e_1,e_2$.
The word $aba\cdots$ lies in $e_1T_de_1$ if $d$ is even and in $e_1T_de_2$ if $d$ is odd; $bab\cdots$ lies in $e_2T_de_2$ ($d$ even), $e_2T_de_1$ ($d$ odd). So
\[
h_{T_S(V)}(t)=\frac1{1-t^2}\begin{pmatrix}1&t\\t&1\end{pmatrix}.
\]
\end{solution}
\subpart Find a basis of $A=T_S(V)/(ab)$ and show that $\dim_\kk A<\infty$.
\begin{solution}
The ideal $(ab)$ is spanned by words containing $ab$; every alternating word of length $\geq3$ contains $ab$ ($aba$, $bab$), so a basis is
\[
  e_1,\ e_2,\ a,\ b,\ ba\qquad(ba\in e_2Ae_2),
\]
and $\dim_\kk A=5$.
\end{solution}
\end{subparts}


\part Let $A$ be a locally finite-dimensional graded algebra generated in degree $1$ over the finite-dimensional semisimple algebra $S=A_0$.
\begin{subparts}
\subpart 
Fix a bimodule $ {}_SV_S$ and an $S$-bimodule isomorphism $\varphi:V\xrightarrow{\sim}A_1$.
Prove that
\[
h_A(t)=h_{T_S(V)}(t)
\quad\Longrightarrow\quad
T_S(V)\xrightarrow{\sim}A, 
\]
where the isomorphism of graded algebras on the right-hand side is induced by $\id_S$ and $\varphi$.

\subpart Consider $\kk[x,z]/(x^2)$ with $\deg x=1$ and $\deg z=2$.  Show that
\[
h_B(t)=h_{T_\kk(B_1)}(t)
\qquad\text{but}\qquad B\not\cong T_\kk(B_1)
\]
as graded algebras.
\end{subparts}

\begin{solution}
For (a), the universal property gives a surjective graded algebra map
$\mu:T_S(V)\to A$ with $\mu|_S=\id_S$ and $\mu|_V=\varphi$.
Put $K=\ker\mu$, and choose primitive idempotents $e_1,\ldots,e_n\in S$ so that the $e_iS$ represent the simple right $S$-modules once each.
For every $d\geq0$ and $1\leq i\leq n$,
\[
0\longrightarrow e_iK_d\longrightarrow e_iT_S(V)_d
\longrightarrow e_iA_d\longrightarrow0
\]
is exact as a sequence of right $S$-modules.
Equality of Hilbert series gives $[e_iK_d:e_jS]=0$ for every $1\leq j\leq n$, hence $e_iK_d=0$.
Since the $e_i$ include a primitive idempotent in each matrix factor of $S$, Artin--Wedderburn gives $S(\sum_{i=1}^n e_i)S=S$.
Therefore $K_d=0$ for every $d$, and $\mu$ is an isomorphism.

For (b), let $B=\kk[x,z]/(x^2)$ with the stated degrees, writing $x,z$ also for their residue classes.
The elements $\{z^q,xz^q\mid q\geq0\}$ form a basis, so
\[
h_B(t)=\sum_{q\geq0}(t^{2q}+t^{2q+1})
=\frac{1}{1-t}=h_{T_\kk(B_1)}(t).
\]
Here $B_1=\kk x$, but $z$ does not belong to the subalgebra generated by $x$.
Also $x\neq0$ and $x^2=0$, whereas non-zero polynomials in $T_\kk(B_1)\cong\kk[t]$ have non-zero squares.
\end{solution}

\end{parts}

\end{questions}

\end{document}
